\begin{figure}[!tbp]
	\centering
	\begin{tikzpicture}[
		x=1cm,y=1cm,
		particle/.style={draw=black!75,thick,fill=blue!16},
		tested/.style={draw=blue!70!black,very thick,fill=blue!28},
		clear ray/.style={-{Latex[length=2.2mm]},very thick,green!55!black},
		blocked ray/.style={-{Latex[length=2.2mm]},very thick,red!75!black},
		label text/.style={font=\small,align=center}
	]
		% Neighbouring particles are drawn first so that the rays remain visible.
		\foreach \x/\y/\r in {
			4.55/0.55/0.63,
			4.25/2.65/0.72,
			2.30/3.45/0.67,
			0.45/2.75/0.62,
			0.30/0.60/0.68,
			2.15/-1.15/0.63}
			\draw[particle] (\x,\y) circle (\r);

		% Particle whose exposure is evaluated.
		\coordinate (C) at (2.15,1.25);
		\draw[tested] (C) circle (0.78);
		\fill[blue!70!black] (C) circle (1.5pt);
		\node[label text,below right=6.5mm and -7mm of C]
			{particle $i$\\centre $\boldsymbol{x}_i$};

		% Unobstructed rays extend into the surroundings.
		\draw[clear ray] (2.6,1.93) -- (3.70,3.50);
		\draw[clear ray] (1.40,1.45) -- (-0.50,1.96);
		\draw[clear ray] (1.58,0.70) -- (-0.20,-1.00);

		% Blocked rays stop at the first neighbouring particle.
		\draw[blocked ray] (2.90,1.03) -- (3.95,0.73);
		\draw[blocked ray] (2.80,1.68) -- (3.65,2.25);
		\draw[blocked ray] (1.56,1.77) -- (0.92,2.34);

		% Ray-start offset and direction annotation.
		\fill[black] (2.6,1.93) circle (1.4pt);
		\node[label text,above right=1mm and 10mm of C,shift={(0,0)}]
			{$\boldsymbol{x}_i+(r_i+\varepsilon)\boldsymbol{s}_m$};
		\node[label text,right] at (2.95,3.5) {$\boldsymbol{s}_m$};

		% Compact legend.
		\draw[clear ray] (6.35,2.65) -- (7.25,2.65);
		\node[label text,right] at (7.28,2.65)
			{unobstructed\\$b_{im}=0$};
		\draw[blocked ray] (6.35,1.35) -- (7.25,1.35);
		\node[label text,right] at (7.28,1.35)
			{blocked\\$b_{im}=1$};
	\end{tikzpicture}
	\caption{Schematic of the ray-tracing calculation for one particle $i$.
		Rays start just outside its surface and follow sampled outward directions
		$\boldsymbol{s}_m$. A ray is blocked when it intersects another particle;
		otherwise, it contributes to the exposed fraction. For clarity, only a
		two-dimensional section and a small subset of rays are shown; the actual
		calculation distributes $M=512$ directions over the complete sphere.}
	\label{fig:ray_tracing_exposure}
\end{figure}
